\section*{Chapitre \ref{ch:fm:limits-and-capital-allocation} --- Limites et allocation du capital}

\begin{solution}{exo:fm:limits-and-capital-allocation:1}
$7.6/8=95\%$ de la \omterm{def:fm:limits-and-capital-allocation:breach}{limite ferme}, au-dessus du seuil indicatif de 6~: un \omterm{def:fm:limits-and-capital-allocation:breach}{dépassement de limite} indicative. Avec une \omterm{def:fm:limits-and-capital-allocation:breach}{limite ferme} à 12~: 63~\%, toujours au-dessus du seuil indicatif.
\end{solution}

\begin{solution}{exo:fm:limits-and-capital-allocation:2}
Pris isolément~: $10.9/65.5=16.7~\%$, valeur ajoutée $10.9-0.15\times65.5=1.12$. Euler~: $10.9/48.0=22.8~\%$, valeur ajoutée 3.74.
\end{solution}

\begin{solution}{exo:fm:limits-and-capital-allocation:3}
$393.9-212.8=181.1$, soit 46.0~\% de la somme des capitaux pris isolément.
\end{solution}

\begin{solution}{exo:fm:limits-and-capital-allocation:4}
Au RAROC d'Euler~: 11.1~\% pour la vente de volatilité, 12.6~\% pour le carry de crédit.
\end{solution}

\begin{solution}{exo:fm:limits-and-capital-allocation:5}
Il gagne dans les années de krach qui forment la queue de la firme, donc son P\&L espéré y est positif et sa contribution négative. Un desk auquel on facture un
montant négatif reçoit un crédit pour avoir ajouté du risque~; mais la contribution varie à mesure qu'il grossit (un grand livre de suivi de tendance domine la queue elle-même), donc le
crédit doit être plafonné.
\end{solution}

\begin{solution}{exo:fm:limits-and-capital-allocation:6}
En ajoutant~: $6\times0.7\times8=33.6>30$, un \omterm{def:fm:limits-and-capital-allocation:breach}{dépassement de limite}. Avec une corrélation de 0.2 et une valeur en risque proportionnelle à la volatilité~: $5.6\sqrt{6+30\times0.2}=19.4$, pas de dépassement.
La limite de la firme doit être vérifiée sur la mesure qui la définit.
\end{solution}

\begin{solution}{exo:fm:limits-and-capital-allocation:7}
Suivi de tendance 24.7 et vente de volatilité 20.3 (ES de la firme~: 114.7)~: sans le krach, le suivi de tendance est une stratégie ordinaire et volatile, et la vente de volatilité paraît bon marché.
L'allocation ne vaut que ce que valent les queues des scénarios.
\end{solution}

\begin{solution}{exo:fm:limits-and-capital-allocation:8}
Les limites par desk ne reconstituent pas l'\omterm{def:fm:the-head-of-desks-job:appetite}{appétit pour le risque} de la firme, sauf si le dispositif les y amène~: des desks à 70~\% peuvent cumuler un dépassement au niveau de la firme, et les risques qu'aucun desk
(concentration entre les desks, krach commun) n'apparaissent qu'au niveau de la firme.
\end{solution}

\begin{solution}{pb:fm:limits-and-capital-allocation:1}
\begin{enumerate}
\item Voir les \cref{def:fm:limits-and-capital-allocation:framework,def:fm:limits-and-capital-allocation:breach}.
\item 95~\% de la \omterm{def:fm:limits-and-capital-allocation:breach}{limite ferme}, un \omterm{def:fm:limits-and-capital-allocation:breach}{dépassement de limite} indicative~; 63~\% après la hausse, toujours un \omterm{def:fm:limits-and-capital-allocation:breach}{dépassement de limite} indicative.
\item 72~\% (21.5~millions de dollars contre 30~millions de dollars).
\item Les desks peuvent chacun rester dans leur limite alors que leur somme dépasse celle de la firme, et un desk peut dépasser la sienne alors que la firme reste largement en deçà.
\item La reconduire à chaque échéance sans l'approbation qu'exige un changement permanent.
\item L'expected shortfall (ES) à 99~\% du P\&L annuel~: 212.8~millions de dollars~; P\&L espéré de 59.1~millions de dollars.
\item Pris isolément 65.5, 69.0, 76.9, 122.8, 47.9, 11.8~; Euler 48.0, $-21.5$, 64.0, 108.5, 13.9, $-0.1$~; incrémentale 42.5, $-31.2$, 59.6, 93.6, 10.0,
$-0.7$ (millions de dollars).
\item Voir la \cref{prop:fm:limits-and-capital-allocation:euler}~; la colonne d'Euler totalise 212.8.
\item L'expected shortfall est sous-additif~: la diversification signifie que les queues ne surviennent pas toutes ensemble.
\item Le suivi de tendance et le market making, dont la contribution est presque nulle~: il gagne dans les pires années de la firme.
\item Voir la \cref{def:fm:limits-and-capital-allocation:raroc}.
\item Pris isolément~: carry de crédit, vente de volatilité, macro. Euler~: carry de crédit et vente de volatilité.
\item Pris isolément 1.12, 1.66, $-3.50$, $-6.34$, $-1.17$, 8.24~; Euler 3.74, 15.24, $-1.55$, $-4.20$, 3.93, 10.02.
\item Un RAROC indéfini ou négatif, et un crédit qui récompense une croissance sans limite~; les firmes plafonnent les crédits ou appliquent un plancher.
\item L'allocation incrémentale, qui mesure le capital de la firme avec et sans la stratégie.
\item 2, 2, 1.57, 0.57, 2, 2~; valeur ajoutée de 27.2 à 59.0~millions de dollars par an.
\item La direction (réduire les stratégies à queue épaisse, faire croître les autres), et non la taille, que fixe le plafond à 2.
\item Oui, à condition qu'elle ait une échéance, que le taux d'utilisation de la firme le permette (72~\%) et que la contribution du desk à la queue de la firme, à la taille supérieure, soit facturée~;
son RAROC d'Euler (11.1~\%) est déjà inférieur au taux de rendement minimal exigé, donc un portefeuille plus grand détruit davantage de valeur.
\item Carry de crédit~: 10.5~\% pris isolément et 12.6~\% en Euler~; vente de volatilité~: 9.8~\% et 11.1~\%~; suivi de tendance~: 69.0~millions de dollars pris isolément, $-21.5$~millions de dollars en Euler.
\item Répartir la \omterm{def:fm:limits-and-capital-allocation:raroc}{charge en capital} selon les contributions d'Euler avec un plancher, réduire la vente de volatilité et le carry de crédit, qui ne rémunèrent pas leur part de la queue~;
et n'accorder la hausse temporaire qu'avec une échéance et la charge appliquée.
\end{enumerate}
\end{solution}

\begin{solution}{iq:fm:limits-and-capital-allocation:1}
Une \omterm{def:fm:limits-and-capital-allocation:breach}{limite indicative} alerte et déclenche un réexamen (une valeur en risque à 75~\% de sa limite)~; une \omterm{def:fm:limits-and-capital-allocation:breach}{limite ferme} doit être corrigée ou remontée immédiatement (une limite de perte qui arrête
le trading).
\iqlookfor{la différence d'action requise.}
\end{solution}

\begin{solution}{iq:fm:limits-and-capital-allocation:2}
Chaque mesure laisse échapper quelque chose~: la valeur en risque dépend de son modèle et de son historique, les sensibilités des petits mouvements, les notionnels ne disent rien~; les limites de stress, de perte et
de concentration couvrent ce qu'elles laissent échapper.
\iqlookfor{angles morts complémentaires.}
\end{solution}

\begin{solution}{iq:fm:limits-and-capital-allocation:3}
Par les contributions d'Euler, la perte espérée de chaque stratégie dans la queue de la firme~; elles totalisent 90. Cinquante chacune ignore la diversification et compte deux fois.
\iqlookfor{additivité à la mesure de la firme.}
\end{solution}

\begin{solution}{iq:fm:limits-and-capital-allocation:4}
$12/108=11.1~\%<15~\%$~: il détruit 4.2 par an. Un autre modèle de queue, ou une valeur de franchise hors de son P\&L, pourrait le modifier.
\iqlookfor{RAROC face au taux de rendement minimal exigé, et le modèle derrière le capital.}
\end{solution}

\begin{solution}{iq:fm:limits-and-capital-allocation:5}
Elle gagne de l'argent dans les pires scénarios de la firme. Non~: sa contribution est marginale à sa taille actuelle et devient positive dès qu'elle domine le portefeuille.
\iqlookfor{marginale, et non moyenne.}
\end{solution}

\begin{solution}{iq:fm:limits-and-capital-allocation:6}
Ajouter à la simulation des scénarios de krach, historiques et hypothétiques (livre~6, chapitre~22), avec la perte de chaque stratégie dans ces scénarios estimée à partir de
ses positions, puis allouer le capital sur l'ensemble combiné des scénarios.
\iqlookfor{scénarios au-delà de l'échantillon.}
\end{solution}
